\documentclass{article}\usepackage[T1]{fontenc}\begin{document}
For fixed split endpoint data, NativeBlocks constructs a functor from native strings of native blocks to the previous single coproduct of span powers. Both source and destination commute as equalities of functors, so native moving-endpoint arrows are included. Flattening a single outer block is the identity functor on the whole tagged span. Replacing every primitive grade by a singleton block and then flattening is also the identity functor. These are component laws on this explicit presentation, not a packaged monad on a category of spans.

For fixed blocks, the original base and grade categories share object universe b and hom universe w as in CoproductSpan.All.span, with no size relation between b and w. For all nested words at once, the base and grades have object universe u and hom universe max(u,v); the latter is closed under the hom encoding of the sigma category used for length tags. There is no payload presheaf in this construction. Split endpoint laws construct the native destination functors. Choice is used for a strict inverse to a fully faithful functor bijective on objects and for equivalences to iso-comma categories, not to redefine the native powers. Normal cleavages are not assumed.

Read lines 929--953 of projects/contextual-composition/foundations/relative-construction/delivery/One\_relative\_construction\_for\_context\_content\_and\_effects.tex at 458ff36a4798e432dee5947afc402f9903a47f14. Here the length-one span unit and multiplication of strings of strings concern the input spans, before application of the presheaf constructor. Read SpanApex.strictFactor and positivePowers, then CoproductSpan.All.span and SpanPull.All.parts. Do not start by substituting the zero-unit PackedFree.monad or the discrete FiniteTelescope.flatten. Check against PR 206 and its frozen source archive before treating the new native construction as related by a monad comparison.

First strengthen the existing strictFactor: the identity-matched last-grade pullback has an object inverse, hence is strictly isomorphic to the Grothendieck snoc apex. Cut.snoc pastes this last-grade pullback with a cut of the prefix. The two reconstruction equations and uniqueness for matched parameter diagrams induce Cut.strictIso: an actual (m+n)-power and the strict pullback of the m- and n-powers are isomorphic categories. Both outer endpoint functors are preserved exactly, and rejoin is a functor equality. For an unrestricted boundary isomorphism the conclusion changes: inclusionEquivalence absorbs it by a cartesian lift in the second span, and equipment gives an equivalence, not a strict inverse to arbitrary iso-comma matching.

The variable tags are part of the grade, not of a terminal payload. splitPairs and gatherPairs are strictly inverse functors: a native arrow has to keep both length tags and still contains both actual block arrows and their matching square. binary concatenates such blocks using Cut.joinFunctor at their unique pair of tags. Nested.word is the strict sequential product of the actual tagged span, with identity span for no block and the tagged span itself for one. At every further step wordEquipment compares strict matching to the corresponding iso-comma product via the whole tagged source cleavage. Nested.flatten sums the collapse functors; flatten\_length computes the total of the inner lengths including zeros. flatten\_single and singletons\_flatten are equalities on all morphisms. For the second law, singletonBlocks cuts off the last native primitive, recursively inserts singleton blocks into the prefix, and rejoin recovers the entire uncut morphism.

Moving.lean compares an actual native path of two blocks (lengths 1 and 2) over natural-number contexts against each of the three raw snoc positions, without defining those projections through the inverse cut. The ending context moves from 11 to 16 along a noninvertible arrow, and the final target grades are 3,4,7 in that order. A ragged word with inner lengths 2,0,1 flattens to length 3 and ends at 11. Parallel.lean uses one-object categories of native grades with word-valued morphisms: the two projected positions of the flattened composite have words [1,2,5] and [8,9,3]; reversing the maps gives [9,3,8] in the second position and a different arrow. IsoBoundary.lean gives the Boolean swap as a boundary isomorphism between two empty powers over Type; it proves it is in the essential image of the comparison but is not literally an image object. This is the reason for reporting equivalence rather than isomorphism on the iso-comma side.

The result is a span map of the concrete strict-match presentations at fixed split endpoints. Naturality in maps of spans, preservation by this map of cartesian arrows or chosen normal lifts, and associativity on three nesting levels are not established. Successive iso-comma products have the stated equivalences; no strict law for arbitrary normalization choices or coherence of all reassociation isomorphisms follows here. Compatibility with the packed free-monad multiplication, the joint base/payload comparison, and the equivalence of strict telescope algebras with strict contextads remain open. Issue \#105 stays open.
\end{document}
